De jager nam de tijd
en bekeek toegewijd

hetgeen was gevangen.
Door de kou bevangen

begon zij te rillen
als haar kleine billen

zachtjes werden bemind
door een licht vlaagje wind.

De koude winter zond
uit de bevroren grond

zijn ijzige groeten
naar haar kleine voeten.

Het weerloze zwerfkind,
gehard door weer en wind,

kon de kilte laten.
Het afzien kon baten.

Als ze hem liet begaan
dan zou hij weer weggaan.

“Mijn tere vrijheid zal
net als de gouden bal

aan de schone prinses
weer na een levensles

worden teruggebracht.”
De kikker had de macht

en was hier als jager
de nare aanjager

van Emmy’s avontuur.
Voor het zoet eerst het zuur.

De jager Actaeon
vond bij een rivierbron

ook een naakte schone,
maar zij bleek geen gewone

badende nimf te zijn
maar uit Zeus’ bloedlijn.

Het beeld werd toegepast
in een perspectiefkast.

De jager kreeg de doos
als een van de cadeaus

waarmee pa hem afkocht
als hij de vrouw opzocht

die zijn minnares was.
Vader was mecenas

en uiterst welgesteld.
Als dank voor zijn giftgeld

kreeg hij deze kijkdoos.
Zijn zoon keek ademloos

dagelijks naar de show.
De jager leerde zo

over Griekse mythen,
maar kon ook genieten

van de blote dozen
die zo moesten pozen.

Toen de jager vernam
dat de scène voortkwam

uit een groter verhaal
ging ermee aan de haal

zijn rijke verbeelding.
Alsof de afbeelding

met zijn aanwijsborden
levend was geworden.

De godin Artemis
ondervond ergernis

van de grove inbreuk.
Ze sprak een toverspreuk

en Actaeon werd
onvoorbereid tot hert,

die door zijn jachthonden
levend werd verslonden.

De jager leerde vroeg
dat hoe men zich gedroeg

een God boos kon maken.
Zelfs als de afspraken

de mens onbekend zijn.
Dat feit gaf hem hoofdpijn.

Een goede intentie
is dus geen garantie

op een veilig leven.
De jongen moest beven

van die onzekerheid
en trieste kwetsbaarheid.

Die kijk op de kijkdoos
maakte hem hulpeloos.

Tot ontwijken genoopt
werd de gluurkist gesloopt,

als het angst-agressief gedrag
dat men vaak terugzag.

De jager pakte haar
bij haar middel zowaar

en boog door zijn knieën
om zijn fantasieën,

over wat er misschien
in haar pruim was te zien,

tegemoet te komen.
Emmy liet hem dromen

door zich weg te draaien
en snel af te taaien.

De oude man bleek rap
en gaf haar toen een klap.

“Wil je dat ik je dood?”
Die bevoegdheid genoot

iedere nabije
bij een vogelvrije.

Ze wenste hem niet boos
en toonde haar kijkdoos.

“Wat zocht hij in haar gat?
Had haar kijkkast iets dat

hem kon amuseren?”
Emmy zou snel leren

wat deze man daar vond.
Voor nu hield zij haar mond.

De jager voelde angst
niet alleen bij zijn vangst,

maar tevens in zichzelf.
Nog wetend het gewelf

van een kelder waarin
toendertijd zijn vriendin

naakt op een tafel lag.
Zijn liefje had ontzag

voor een sekteleider,
die als begeleider

van haar en de jager
hen bijstond om trager

bij seks te werk te gaan.
Hoe ze beiden voldaan

zonder klaar te komen
hun vuur lieten stromen.

Spiritualiteit
als excuus voor geilheid.

Een paar sekteleden
waren aangetreden.

Ze mochten niet praten,
maar wel door kijkgaten

in de stenen muren
naamloos er naar gluren.

Wel naar de show piepen,
maar zonder te piepen.

De jager vond dat taai.
“Als alleen ik je naai

en anderen kijken
kan het ons verrijken.”,

loog hij stoer tegen haar.
Zijn uiting bleek niet waar.

Angst haar te verliezen
deed hem ervoor kiezen

om steeds te doen alsof.
Hij vond de leer niet tof,

maar het bracht hen samen.
Dus bleef hij beamen

uit zijn verlatingsangst.
Oneerlijk duurt het langst.

“Niet klaarkomen! Hou vol!”
De jager kon zijn rol

echter niet waarmaken
en moest haar volmaken.

“Bij Jezus, wat een oen!
Ik zal het wel voordoen!”

De jager keek verschrikt,
maar had eerder geslikt

de pil van overspel.
Hij kende het dus wel.

Na slechts zestig tellen
hoorde men vertellen:

“Ook ik ben maar een mens”,
zuchtte de coach intens.

De jagers inspectie
mistte introspectie.

Kwam de angst uit dreiging?
Of was het een zinspeling

vanuit eerdere angst.
De jager werd het bangst

van innerlijke vrees,
die naar oud zeer verwees.

Angst als herinnering
aan een teleurstelling

voorspelt dit als uitkomst
opnieuw in de toekomst

en zaait het bovendien.
Kon de jager dat zien?

Angst als vriend omarmen?
Hij had geen erbarmen

aangeleerd voor zwakte.
Het woord dat men plakte

op slappe emoties.
Het waren erupties

om je voor te schamen
als ze binnenkwamen.

Zette al zijn zinnen
op het overwinnen.

Blijft angst wel aanwezig
als we niet meer bezig

zijn het te bevechten?
Moeten we het knechten?

Wie is dan de vechter?
Wie is dan de knechter?

Of ben je zelf de angst?
Maakt verzet ons het bangst?

Is het overgeven
en intens beleven

een betere aanpak?
Wordt diens grip dan niet zwak?

Deze vragen stelde
de jager maar zelden.

Hij stelde wel de vraag;
"Wat doe ik nog meer graag?”
